% MAS1004 과제 1: 짧은 보고서, 한국어 양식.
%
% 이 과제를 하며 배운 것을 자유롭게 작성하면 됩니다.
%   - 어떤 주제를 골랐고, 왜 골랐는지
%   - 과제가 어떻게 진행되었는지
%   - 무엇을 배웠는지
%
%  이 리포트는 본문을 한쪽으로 작성하며, LLM을 쓰지 않고 직접 작성합니다.
% 본문 뒤에 그림을 세 쪽까지 더 붙일 수 있습니다.
% XeLaTeX로 컴파일하면 됩니다.  (Overleaf: Menu, Compiler, XeLaTeX).

\documentclass[11pt]{article}
\usepackage[onepage,ko]{mas1004}

\reporttitle{과제 1: 짧은 보고서}
\studentname{김서영}
\studentid{20251120}

\begin{document}
\makeheader

% 본문, 한 쪽. 아래 안내에 필요한 LaTeX 사용법이 있습니다.
% 읽은 뒤 지우고, 그 자리에 보고서를 작성하면 됩니다.

나는 평소 서브컬쳐 게임을, 특히 다양한 매력을 지닌 캐릭터의 디자인과 서사가 있는 명조 (Wuthering Waves) 라는 게임을 가장 즐겨한다. 이 서브컬쳐 장르에서는 공식 일러스트뿐만 아니라 팬아트 등 다양한 창작물이 만들어지는데, 이러한 그림들을 찾아 보는 걸 굉장히 좋아하는 편이다. 한 캐릭터를 좋아하는 사람들이 매우 많다 보니, 같은 캐릭터도 창작자에 따라 서로 다른 그림체와 방식으로 표현된다. 예를 들어 한 사람은 캐릭터를 SD 그림체로 그리는가 하면, 어떤 사람은 LD 그림체로 그린다. 그 캐릭터를 모르는 사람이 본다면 같은 캐릭터라고 생각할 수 없을 정도이다. 그러나 나는 이러한 차이에도 불구하고 그 캐릭터를 좋아하는 사람들이 같은 캐릭터임을 알아보는 이유와, 그 과정에서 어떤 시각적 특징을 단서로 삼는지 궁금해졌다.

명조 캐릭터 중에 나는 카멜리아, 에이메스, 양양, 린네를 골랐다. 카멜리아는 밈으로 많이 활용돼 다양한 그림체가 있을 것이라고 생각했고, 변형된 이미지도 잘 인식하는지 궁금했다. 양양은 명조의 마스코트이자 앱 아이콘까지 장식하고 있기에 의미가 깊었고, 외형이 달라진 버전도 최근 출시되어 같은 캐릭터로 인식하는지 궁금했다. 에이메스와 린네는 개인적으로 좋아하는 캐릭터라서 골랐다.

학습 전에는 모델이 주로 머리카락 색, 헤어스타일, 의상 등을 기준으로 캐릭터를 분류할 것이라고 예상했다. 특히 팬아트에서는 캐릭터가 원래의 의상과 다른 옷을 입은 모습으로 그려지는 경우가 많기 때문에, 의상보다는 머리카락 색과 헤어스타일이 분류에 더 큰 영향을 미칠 것으로 생각했다.

먼저 사진을 수집하는 작업에는 인공지능을 활용했다. 이후 수집된 자료를 정리하고, 학습에 사용할 사진의 선정 기준을 정해 선별했다. 일단 1차적으로 모은 사진 중 잘못 가져온 사진을 걸러내는 건 인공지능이 진행하고, 2차적인 과정은 내가 직접 진행했다. 이 때 내가 정한 사진 선별 기준은 캐릭터가 이미지의 50\% 이상을 차지하고, 머리카락이나 의상의 특징이 어떤 식으로든 식별 가능한지를 중요하게 보았다. 학습 과정에서는 학습 횟수를 10 epochs 로 유지하면서 학습률을 0.0001 에서  0.0003 으로 변경하는 실험을 진행했다. 학습 횟수를 증가하면 모델이 학습하면 할수록 이상한 결과를 내놓는 일이 내 경험상 비일비재했어서, 학습 횟수는 고정하면서 학습률 변화의 영향을 살펴보고자 한 선택이었다.

완성된 모델에 직접 사진을 입력하여 평가한 결과, 원본 개인 사진 24중 22장을 올바르게 분류해 91.7\%의 정확도를 기록했다. 다만 사람이 보기에는 분명히 '린네'와 '카멜리아'인 두 사진을 모델이 모두 '에이메스'로 분류한 점이 눈에 띄었다. 두 사진을 공통점을 살펴보니, 사진 선별의 기준이었던 캐릭터가 이미지의 50\% 이상을 차지하지 않았다. 그래서 캐릭터 중심으로 크롭한 뒤 결과를 비교하는 별도 실험을 진행했다. 크롭 후 캐릭터에 대한 예측 점수는 카멜리아가 0.9\%에서 100\%로, 린네가 17.2\%에서 67.5\%로 상승했으며, 두 사진 모두 올바른 캐릭터로 분류되었다.

그러나 크롭이 항상 도움이 된 것은 아니었다. 정확도가 비교적 떨어졌던 하위 4개의 사진을 다시 크롭하여 재측정도 해보았는데, 다른 '린네' 사진에서는 예측 점수가 80.9\%에서 68.5\%로 낮아졌다. 따라서 이 실험을 통해 일부 사진에서 크롭이 분류를 개선할 수 있음을 보여주지만, 모든 사진에 같은 효과가 나타난다고 보기는 어렵다. 또한 크롭은 배경뿐 아니라 캐릭터의 크기와 구도도 함께 바꾸므로, 점수 변화만으로 어떤 요소가 결과에 영향을 주었는지 확정할 수는 없다. 다만 내 예상으로 위의 '린네' 사진의 예측도가 낮아진 이유는, 본래 정사각형 모양에 가까웠던 사진을 직사각형 모양으로 캐릭터만 정확히 잘라내어서, 모델이 정사각형으로 다시 잘라서 측정할 때 '린네'로 분류할 수 있는 특징이 잘려나갔던 게 아닐까 싶다.

실험 전에는 비교적 윤곽이 뚜렷한 인형탈 사진을 쉽게 맞힐 것이라고 예상했다. 하지만 사람이 쉽게 알아볼 수 있는 사진이라고 해서 모델도 반드시 올바르게 분류하는 것은 아니고, 또 반대로 사람이 쉽게 알아볼 수 없는 사진이라고 해서 모델이 반드시 틀리게 예측을 하는 것은 아니라는 점을 알게 되었다. 이번 실험을 통해 입력 이미지의 구성을 바꾸어 결과를 비교하는 것이 유용할 수 있으며, 예측 점수가 달라졌다는 사실과 그 변화의 원인을 구분해 해석해야 한다는 점을 배웠다.

더불어, 앞서 예측했던 '헤어스타일이 분류에 더 큰 영향을 미칠 것이다'라는 것에 대해서는 에이전트가 이미지를 편집하고 모델을 분류했다. 머리카락을 가리자 정답이 22장에서 20장, 실제 캐릭터 점수가 19장에서 하락했다. 카멜리아 분류는 다만 그대로였기 때문에, 옷 같은 다른 단서로 이용했을 수 있다고 생각했다. 간단히 배운 점이라함은 머리카락이 영향을 준다는 예상은 뒷받침되었지만, 가장 큰 영향을 줬는지는 미지수이다.

\figurepages

% 그림, 세 쪽까지.

% \begin{figure}[htbp]
%   \centering
%   \includegraphics[width=0.7\linewidth]{example-image-a}
%   \caption{그림 설명.}
% \end{figure}

\end{document}